% ICML SCALE 2026 Workshop Submission - Double-Blind Review
\documentclass[10pt,twocolumn,letterpaper]{article}
\usepackage{icml2026}

% Math and algorithm packages
\usepackage{amsmath,amsfonts,amssymb}
\usepackage{algorithm}
\usepackage{algorithmic}
\usepackage{graphicx}
\usepackage{booktabs}
\usepackage{url}
\usepackage{xcolor}
\usepackage{microtype}

\icmltitlerunning{IdleKV: Exploiting Agentic GPU Idle Time for KV Cache Recovery}

\begin{document}

\icmltitle{IdleKV: Exploiting Agentic GPU Idle Time for KV Cache Recovery}

\icmlsetsymbol{equal}{*}

\begin{icmlauthorlist}
\icmlauthor{Anonymous Authors\icmlequal}{anon}
\end{icmlauthorlist}

\icmlaffiliation{anon}{Anonymous Institution}

\icmlcorrespondingauthor{Anonymous Authors}{anon@domain.com}

\maketitle

% ---------------------------------------------------------------
%  Abstract
% ---------------------------------------------------------------
%
% CONVENTIONS APPLIED:
%  - 5-sentence workshop formula: problem → gap → method → mechanism → result+significance
%  - No hedging on core claims; numbers bound every result claim
%  - "We propose" in sentence 3; method name bolded
%  - Significance woven into final sentence
% ---------------------------------------------------------------
\begin{abstract}
Agentic LLM systems waste 30--55\% of GPU time during tool calls
\citep{continuum2025,infercept2025}, yet KV cache compression methods make
permanent eviction decisions during generation, when both information and
compute are scarce.
Existing quality-recovery methods such as RefreshKV \citep{refreshkv2025}
recover up to 52\% of the quality gap but impose periodic latency on every
decode step, degrading throughput by ${\sim}25\%$.
We propose \textbf{IdleKV}, an idle-time cache-repair framework for
tool-calling LLMs.
Its core short-idle operator performs TIR-informed token re-scoring over a
bounded shadow buffer in under 100\,ms, after compression has already happened
and after the effective query has had a chance to shift.
On a delayed-query stress gate at $r{=}0.7$, the compressed baseline improves
from 91.7\% to 100.0\% with a mean 27.4\,ms Phase~1 pass.
On the informative Qwen2.5-7B RULER-4K slice, 100\,ms Phase~1 improves average
accuracy from 0.877 to 0.903 across follow-up seeds.
Longer-budget full-attention refresh remains exploratory; the current measured
scope supports the narrower claim that short-idle repair can turn otherwise
wasted tool-call time into usable inference-quality budget.
\end{abstract}

% ---------------------------------------------------------------
%  1. Introduction
% ---------------------------------------------------------------
%
% CONVENTIONS APPLIED:
%  - Opens with trend + specific numbers in first two sentences
%  - Gap paragraph names methods and precise failure modes; cites 7 works
%  - "Our approach" paragraph gives technical intuition with at least one number
%  - 3 contribution bullets: present tense, bold labels, number in every bullet
%  - Related work merged into intro gap paragraphs (workshop convention)
%  - No "rest of the paper is organized as follows" paragraph
% ---------------------------------------------------------------
\section{Introduction}
\label{sec:intro}

Modern agentic LLM deployments---Cursor, Claude Code, SWE-agent---spend
30--55\% of wall-clock session time with tool calls in flight.
Continuum \citep{continuum2025} documents up to 58.2\% of end-to-end latency
consumed by scheduling bubbles and KV-cache recomputation around tool-call
boundaries; InferCept \citep{infercept2025} identifies 37--40\% of forward
time spent recomputing already-processed context under multi-tenant
batching.
In \emph{single-agent} deployments with a dedicated GPU---local coding agents,
on-device assistants, reserved-capacity inference---the accelerator sits
truly idle while a tool executes externally.
This is the setting IdleKV targets; multi-tenant serving with iteration-level
scheduling (vLLM, SGLang) is out of scope.

KV cache compression is necessary for long agentic sessions: without it,
memory pressure forces truncation or eviction of critical context.
However, existing compression methods---H2O \citep{h2o2023},
SnapKV \citep{snapkv2024}, StreamingLLM \citep{streamingllm2023}---make
permanent eviction decisions during generation, when both information and
compute are scarce.
Three findings establish why this is fundamentally suboptimal.
\citet{lazyeviction2025} show that $>$95\% of evicted tokens exhibit
Token Importance Recurrence (TIR) on multi-step reasoning traces (GSM8K,
MATH500): tokens unimportant at eviction often become critical as context
evolves.
The TIR phenomenon has not yet been measured directly on agentic/tool-calling
traces; we treat this figure as a motivating upper bound and report
IdleKV's realized recovery on agentic workloads directly.
\citet{defensivekv2026} demonstrate that attention-score importance is
unstable across decode steps, causing systematic errors in one-shot eviction.
Methods that do attempt quality recovery during generation---RefreshKV
\citep{refreshkv2025}, CaliDrop \citep{calidrop2025}, SeerAttn-R
\citep{seerattnr2025}---impose periodic or per-step latency overhead that
degrades decode throughput by 15--30\%, trading one problem for another.
The core difficulty is a scheduling mismatch: quality recovery is most useful
during generation, but that is precisely when latency constraints make it most
costly.

\textbf{IdleKV} resolves this mismatch by scheduling quality recovery during
tool-call pauses, when the GPU is otherwise idle.
Compression happens during generation, where speed is paramount.
Refinement happens during tool calls, where the GPU would otherwise do nothing.
The current measured core is Phase~1: TIR-informed re-scoring of recently
evicted tokens against recent queries in under 100\,ms, promoting tokens whose
importance has grown since eviction.
We also implement a longer-budget Phase~2 extension based on progressive
full-attention refresh using a stored uncompressed cache, but the present
evidence only supports Phase~1 as the durable main claim.
Neither phase adds work to the token-by-token generation critical path.

\textbf{Contributions:}
\begin{itemize}
\item \textbf{Idle-time, delayed-query cache repair.}
  We formulate a setting where KV compression decisions are made before the
  final effective query is known, and where tool-call pauses provide real idle
  compute that can be used to repair those earlier decisions.
\item \textbf{A cheap shadow-buffer repair operator.}
  IdleKV Phase~1 re-scores recently evicted tokens against recent query states
  in under 100\,ms and moves quality recovery off the generation critical path.
\item \textbf{Measured evidence for the short-idle claim.}
  In the current measured scope, the delayed-query gate and the informative
  multi-seed Qwen RULER-4K slice both support the same conclusion:
  short-idle Phase~1 improves over the matching aggressive-compression
  baseline, while longer-budget Phase~2 remains exploratory.
\end{itemize}

A natural question is whether less aggressive compression achieves the same
quality without the complexity of idle-time refinement.
IdleKV directly addresses this by asking a different question:
once aggressive compression has already happened, can a short idle window
recover some of the lost quality without adding per-token decode work?
The current paper answers that narrower question first.
Broader Pareto claims against less aggressive compression and synchronous
refresh remain part of the larger scale-up study rather than the core measured
claim in this version.

% ---------------------------------------------------------------
%  2. Method
% ---------------------------------------------------------------
%
% CONVENTIONS APPLIED (method section 20-30% of paper):
%  - 1-2 paragraph preamble gives roadmap; no separate "Overview" subsection header
%  - Figure 1 introduced early; system diagram anchors the section
%  - Algorithm environments for low-level kernel-adjacent work
%  - Inline quantitative claims with bounding conditions throughout
%  - "We" for all novel contributions
% ---------------------------------------------------------------
\section{Method}
\label{sec:method}

\textbf{IdleKV} wraps any attention-based compression method in four stages.
During prefill, the base compression method runs and the full uncompressed
cache is stored on CPU.
During generation, the model decodes with the compressed cache while a shadow
buffer maintains the 256 most recently evicted KV pairs per layer per head.
At each tool-call boundary, Phases~1 and~2 run until the tool returns or the
compute budget is exhausted.
On resume, IdleKV atomically swaps the refined cache and continues generation.
The CPU backup updates after each tool-result prefill to maintain consistency.
Figure~\ref{fig:architecture} illustrates the full pipeline.

\begin{figure}[t]
\centering
\fbox{\parbox{0.88\columnwidth}{\centering\vspace{2.0cm}%
\small[Figure~1: Architecture diagram ---
Prefill $\to$ Compress $+$ Store Full KV $\to$
Generate with Shadow Buffer $\to$ Tool Call $\to$
Phase~1 $+$ Phase~2 $\to$ Resume]\vspace{2.0cm}}}
\caption{IdleKV pipeline.  Compression runs during generation
(latency-sensitive path); refinement runs during tool calls (free compute).
The CPU backup is updated after each tool-result prefill.}
\label{fig:architecture}
\end{figure}

\subsection{Phase 1: TIR-Informed Re-scoring ($\leq$100\,ms)}
\label{sec:phase1}

\citet{lazyeviction2025} establish that $>$95\% of evicted tokens exhibit
Token Importance Recurrence: a token unimportant at eviction time can become
critical as queries evolve.
Phase~1 exploits this by re-evaluating recently evicted tokens against
recent per-layer queries.
For each layer $\ell$ and head $h$, we maintain a shadow buffer $S_\ell$
of the 256 most recently evicted KV pairs.
Algorithm~\ref{alg:phase1} details the procedure.

\begin{algorithm}[t]
\caption{Phase 1: TIR-Informed Token Re-scoring}
\label{alg:phase1}
\begin{algorithmic}[1]
\REQUIRE Retained cache $R_\ell$, shadow buffer $S_\ell$,
         recent per-layer queries $Q_{t-W+1:t}^{(\ell)}$ (window $W{=}32$)
\ENSURE Updated retained cache $R'_\ell$, updated shadow buffer $S'_\ell$
\FOR{each layer $\ell = 1, \ldots, L$}
  \FOR{each head $h$}
    \STATE $K_{\mathrm{all}} \leftarrow [R_\ell^{(h)};\; S_\ell^{(h)}]$
    \STATE $\mathrm{scores} \leftarrow \tfrac{1}{W}
           \sum_{i=t-W+1}^{t} Q_i^{(h)} K_{\mathrm{all}}^\top$
    \STATE $\mathrm{idx} \leftarrow \mathrm{top\text{-}}k(\mathrm{scores})$
    \STATE Promote shadow tokens in $\mathrm{idx}$ to $R'_\ell$;
           demote displaced tokens to $S'_\ell$
  \ENDFOR
\ENDFOR
\end{algorithmic}
\end{algorithm}

\textbf{Implementation detail.}
We store the hidden state feeding each layer for the last $W{=}32$ generated
or newly ingested tokens, then project those states through each layer's own
$q$-projection at re-scoring time. This matches the layer-local query geometry
used during prefill compression and avoids reusing a single final-layer state
across all layers.

\textbf{Complexity.}
$O(d \cdot |S_\ell| \cdot H \cdot L)$ total, with $|S_\ell|{=}256$, $H$
heads, $L$ layers.
Wall-clock cost is hardware-dependent; on our target A10G setup we budget
Phase~1 to stay below 100\,ms across all layers of Llama-3.1-8B.

Phase~1 differs from CaliDrop \citep{calidrop2025} in both mechanism and
scheduling: CaliDrop stores pre-computed attention contributions and calibrates
score estimates during generation (adding per-step overhead), while Phase~1
stores full KV pairs and performs token swaps exclusively during idle time,
adding no per-step cost.

\subsection{Phase 2: Progressive Full-Attention Refresh (100\,ms--2\,s)}
\label{sec:phase2}

For idle budgets exceeding 100\,ms, Phase~2 performs progressive full-attention
refresh using the CPU-stored uncompressed cache.
We process layers shallow-first: errors in early layers propagate through
residual connections \citep{refreshkv2025}, so shallow-layer correction
yields the largest downstream quality benefit per millisecond spent.
Algorithm~\ref{alg:phase2} details the procedure.

\begin{algorithm}[t]
\caption{Phase 2: Progressive Layer-wise Refresh}
\label{alg:phase2}
\begin{algorithmic}[1]
\REQUIRE CPU cache $\{K_{\mathrm{cpu}}^{(\ell)}, V_{\mathrm{cpu}}^{(\ell)}\}$,
         budget $T$, current queries $Q$
\ENSURE Partially or fully refreshed retained cache
\STATE $t_0 \leftarrow \texttt{now()}$
\FOR{each layer $\ell = 1, 2, \ldots, L$}
  \IF{$\texttt{now()} - t_0 > T$}
    \STATE \textbf{break}
           \COMMENT{Valid partial result: layers $1\ldots\ell{-}1$ refreshed}
  \ENDIF
  \STATE Load $K_{\mathrm{cpu}}^{(\ell)}, V_{\mathrm{cpu}}^{(\ell)}$
         via PCIe (${\approx}16$\,MB, ${<}1$\,ms)
  \STATE $K_{\mathrm{full}} \leftarrow
         [K_{\mathrm{cpu}}^{(\ell)};\; K_{\mathrm{gen}}^{(\ell)}]$
  \STATE $A \leftarrow \mathrm{softmax}(Q\,K_{\mathrm{full}}^\top / \sqrt{d})$
  \STATE $R_\ell \leftarrow \mathrm{top\text{-}}k(A)$;\; free CPU tensors
\ENDFOR
\end{algorithmic}
\end{algorithm}

\textbf{Anytime property.}
Interruption at any point yields a valid cache: layers $1,\ldots,\ell{-}1$
are refreshed; layers $\ell,\ldots,L$ retain their Phase~1 state.
Cost per layer: 15--40\,ms on Llama-3.1-8B, enabling 2--5 layer refreshes per
100\,ms budget.
Phase~2 reuses the core refresh mechanism from \citet{refreshkv2025} but moves
it entirely out of the generation critical path.

\subsection{Memory and Latency Overhead}
\label{sec:overhead}

Table~\ref{tab:memory} summarizes IdleKV's memory footprint over the SnapKV
baseline.
Total GPU overhead is 40\,MB (${\approx}0.2\%$ of a 24\,GB A10G).
CPU backup for a 4K context requires ${\approx}$512\,MB.
We verify that decode throughput with and without shadow buffer hooks differs
by ${<}2\%$ (not statistically significant at $n{=}3$ seeds), consistent with
the lightweight per-step hidden-state capture design.
Cache state transitions at tool-call boundaries are absorbed within idle
periods by construction.

\begin{table}[t]
\centering
\caption{Memory overhead of \textbf{IdleKV} over SnapKV baseline.
GPU overhead is ${\approx}0.2\%$ of a 24\,GB A10G.}
\label{tab:memory}
\resizebox{0.98\columnwidth}{!}{%
\begin{tabular}{lcc}
\toprule
Component & GPU (MB) & CPU (MB) \\
\midrule
Shadow buffer ($256{\times}L{\times}H_{\mathrm{kv}}{\times}d$) & 32  & ---  \\
Query buffer  ($32{\times}L{\times}d_{\mathrm{model}}$)        & 8   & ---  \\
Full KV backup (4K ctx)                           & --- & 512  \\
\midrule
Total overhead & \textbf{40} & \textbf{512} \\
\bottomrule
\end{tabular}%
}
\end{table}

% ---------------------------------------------------------------
%  3. Experimental Setup
% ---------------------------------------------------------------
%
% CONVENTIONS APPLIED:
%  - Hardware always specified with memory
%  - Baseline list is numbered, each baseline named with its role
%  - Metrics and benchmarks stated explicitly
%  - Detailed hyperparameters deferred to appendix
% ---------------------------------------------------------------
\section{Experimental Setup}
\label{sec:setup}

\textbf{Models and hardware.}
We evaluate on Llama-3.1-8B-Instruct (primary) and Qwen2.5-7B-Instruct
(cross-model validation), both on a single NVIDIA A10G (24\,GB GDDR6).

\textbf{Agentic trace simulation.}
Tool calls are simulated every 100 generated tokens with log-normal idle
duration (median 1\,s, range 100\,ms--10\,s), calibrated to SWE-Bench agentic
traces from \citet{continuum2025}. The broader planned matrix evaluates idle
budgets of 0--5\,s to characterize the full quality-vs-compute tradeoff; the
current measured A10G scout uses a reduced subset (0, 100\,ms, 1000\,ms) to
choose the larger run.

\textbf{Baselines.}
(1)~\textbf{Full cache}: no compression; quality upper bound.
(2)~\textbf{SnapKV $r{=}0.3$}: 30\% token retention; the ``why not compress
less?'' Pareto reference.
(3)~\textbf{SnapKV $r{=}0.5$}: 50\% token retention; the primary throughput
reference in the current measured study.
(4)~\textbf{H2O $r{=}0.5$}: heavy-hitter oracle at 50\% retention.
(5)~\textbf{StreamingLLM}: attention sink $+$ sliding window; lower-bound
quality reference.
(6)~\textbf{Sync refresh}: SnapKV $r{=}0.5$ with RefreshKV-style periodic
refresh during generation (stride~15); isolates the value of idle-time
scheduling from the quality operations themselves in the larger-GPU scale-up
study. In the single-A10G measured study, the central question is whether
IdleKV improves over SnapKV at the same decode operating point while preserving
throughput.

\textbf{Benchmarks and metrics.}
The measured single-A10G core currently consists of the delayed-query pilot and
a reduced RULER-4K scout used to choose the larger matrix. The stable on-host
operating point for the broader A10G study is still RULER at 4K context; we
also run auxiliary LongBench baseline slices on the same hardware. LongBench
with IdleKV and 8K RULER remain scale-up follow-up experiments after
decode-memory optimization or migration to larger-memory GPUs. We report RULER
accuracy, LongBench average score, KL divergence vs.\ full cache, and decode
throughput (tokens/sec).

% ---------------------------------------------------------------
%  4. Results
% ---------------------------------------------------------------
%
% CONVENTIONS APPLIED:
%  - Subsection arc: Main Results → Budget Curve → Pareto → Ablation
%  - 2-4 sentences of prose per table, explicitly calling out key numbers
%  - Key numbers bolded in tables; × notation throughout
%  - Preliminary framing stated once, cleanly, at section start
%  - "we observe" / "we find" for results; passive only for established facts
%  - Ablation at end as dedicated subsection
% ---------------------------------------------------------------
\section{Results}
\label{sec:results}

This version of the paper separates three evidence tiers.
The measured core is the delayed-query pilot plus the reduced single-A10G scout
described above. The broader single-A10G and A100 matrices remain in progress.
The current paper therefore focuses on what is already measured and uses the
remaining tables/figures only as reserved slots for the workshop-core matrix
that runs next.

\subsection{Main Quality Results}

Table~\ref{tab:main_results} summarizes the measured core evidence that already
shapes the workshop-core matrix.
Two conclusions are durable at this stage.
First, delayed-query drift creates recoverable quality loss under aggressive
compression.
Second, short-idle Phase~1 repair is consistently better than the matching
no-idle compressed baseline on the informative Qwen RULER-4K slice.
By contrast, longer-budget P1$+$P2 remains outside the core claim because the
current Qwen follow-up shows a severe regression rather than an incremental
gain.

\begin{table}[t]
\centering
\caption{Measured core evidence that currently supports the workshop claim.
The first row is the delayed-query mechanism gate; the remaining rows are the
informative Qwen RULER-4K scout slices.}
\label{tab:main_results}
\resizebox{0.98\columnwidth}{!}{%
\begin{tabular}{lcccc}
\toprule
Slice & Baseline & Phase~1 @ 100\,ms & P1 only @ 1000\,ms & P1$+$P2 @ 1000\,ms \\
\midrule
Delayed-query pilot (Llama, $r{=}0.7$) & 91.7\% & 100.0\% & --- & --- \\
Qwen scout, seed 42 ($r{=}0.7$) & 0.86 & 0.90 & 0.90 & 0.667 \\
Qwen follow-up mean, seeds 123/456 ($r{=}0.7$) & 0.877 & 0.903 & 0.903 & 0.653 \\
\bottomrule
\end{tabular}%
}
\end{table}

\subsection{Quality-vs-Idle-Budget Curve}

Table~\ref{tab:idle_budget} and Figure~\ref{fig:budget_curve} characterize
RULER accuracy as a function of idle compute budget on the informative Qwen
slice.
The measured data only supports one strong budget claim today:
the first 100\,ms matters.
Going from 0\,ms to 100\,ms helps; going from 100\,ms to 1000\,ms does not
help further for Phase~1 on this slice; and the current P1$+$P2 path is not
stable enough to include in the workshop-core story.

\begin{table}[t]
\centering
\caption{Measured Qwen RULER-4K budget sweep on the informative $r{=}0.7$
slice. The first 100\,ms carries the durable Phase~1 signal.}
\label{tab:idle_budget}
\resizebox{0.98\columnwidth}{!}{%
\begin{tabular}{lcccc}
\toprule
Setting & 0\,ms & 100\,ms & 1000\,ms & Read \\
\midrule
Phase~1, seed 42 & 0.86 & 0.90 & 0.90 & gain saturates by 100\,ms \\
Phase~1, mean over seeds 123/456 & 0.877 & 0.903 & 0.903 & same pattern \\
Phase~1$+$Phase~2, mean over seeds 123/456 & --- & --- & 0.653 & unstable; not a core claim \\
\bottomrule
\end{tabular}%
}
\end{table}

\begin{figure}[t]
\centering
\fbox{\parbox{0.88\columnwidth}{\centering\vspace{1.6cm}%
\small[Figure~2: Quality-vs-idle-budget curves for
$r \in \{0.3,\,0.5,\,0.7\}$.
X-axis: idle budget (ms, log scale).
Y-axis: RULER 4K improvement over no-refinement baseline.]\vspace{1.6cm}}}
\caption{Quality recovery as a function of idle compute budget.
This figure is reserved for the measured workshop-core plot after
`configs/scale_core.yaml` finishes. The current evidence already supports the
same qualitative message: the first 100\,ms matters, while the current
long-budget P1$+$P2 path is not stable enough for the main story.}
\label{fig:budget_curve}
\end{figure}

\subsection{Pareto Frontier}

Figure~\ref{fig:pareto} plots RULER 4K accuracy vs.\ decode throughput for all
methods.
That throughput figure is the next critical measurement rather than a completed
result. The current repository contains a dedicated
`scripts/throughput_spotcheck.py` script for this purpose.
The workshop-core claim only needs the narrower operating-point check:
showing that short-idle Phase~1 improves on the matching compressed baseline
without materially changing decode speed.

\begin{figure}[t]
\centering
\fbox{\parbox{0.88\columnwidth}{\centering\vspace{1.6cm}%
\small[Figure~3: Quality-throughput Pareto frontier.
Each point is one method.
The workshop-core version only needs the short-idle Phase~1 operating-point
check; broader Pareto claims are deferred until measured.]\vspace{1.6cm}}}
\caption{Reserved throughput-quality plot for the workshop-core matrix.
The only required claim for this version is that short-idle Phase~1 improves
quality over the matching compressed baseline without materially changing the
decode operating point.}
\label{fig:pareto}
\end{figure}

\subsection{Ablation: Phase Contribution and Cross-Model Generalization}

Table~\ref{tab:ablation} isolates Phase~1 and Phase~2 contributions.
The current workshop-core ablation is narrower: it asks whether the shadow
buffer is necessary for the shortlisted Phase~1 path.
That is the role of `configs/scale_mechanism_ablation.yaml`.
We explicitly do not treat Phase~2 as part of the current ablation story,
because the Qwen follow-up indicates that the existing long-budget path is not
yet stable enough for a positive claim.

\begin{table}[t]
\centering
\caption{Reserved mechanism ablation slot for the workshop-core matrix.
The active question is whether the shadow buffer matters for the shortlisted
Phase~1 @ 100\,ms path.}
\label{tab:ablation}
\resizebox{0.98\columnwidth}{!}{%
\begin{tabular}{lccc}
\toprule
Configuration & Benchmark & Budget (ms) & Status \\
\midrule
SnapKV $r{=}0.7$ baseline & Qwen RULER 4K & 0 & reserved for measured run \\
IdleKV Phase~1, buffer $0$ & Qwen RULER 4K & 100 & reserved for measured run \\
IdleKV Phase~1, buffer $256$ & Qwen RULER 4K & 100 & reserved for measured run \\
\bottomrule
\end{tabular}%
}
\end{table}

At present, Qwen2.5-7B is the informative model for method selection while
Llama-3.1-8B at RULER 4K is mostly ceiling. The repo therefore includes a small
Llama hardness probe rather than claiming current cross-model parity.

% ---------------------------------------------------------------
%  5. Limitations
% ---------------------------------------------------------------
\section{Limitations}
\label{sec:limitations}

\textbf{Single-tenant scope.}
IdleKV's premise---that GPU compute is free during a tool call---holds
only when a single request owns the accelerator.
Multi-tenant serving engines such as vLLM and SGLang use iteration-level
continuous batching and fill the bubble with other in-flight requests
\citep{continuum2025}; in those settings the real cost is KV-cache churn
(recomputation, memory pressure), not idle compute, and the appropriate
target is systems like InferCept \citep{infercept2025} rather than IdleKV.
We target single-agent deployments (local coding agents, reserved capacity,
on-device inference) where the dedicated-GPU assumption is accurate.
\textbf{TIR domain transfer.}
Our motivating TIR figure \citep{lazyeviction2025} is measured on reasoning
benchmarks; whether the recurrence rate is similar on agentic traces is an
open empirical question that our evaluation partially answers.
We evaluate on 7--8B parameter models; larger models may exhibit different
per-layer refresh costs and shadow buffer requirements.
Integration with adaptive compression methods that adjust token budgets
dynamically (DynamicKV, Ada-KV) is future work.

% ---------------------------------------------------------------
%  6. Conclusion
% ---------------------------------------------------------------
%
% CONVENTIONS APPLIED:
%  - Restates the core result with exact numbers
%  - States the broader principle in the final sentence
%  - No hedging; no bullet points
% ---------------------------------------------------------------
\section{Conclusion}
\label{sec:conclusion}

\textbf{IdleKV} is the first system to exploit GPU idle time during agentic
tool calls for KV cache quality improvement.
The measured core evidence supports a narrower but still useful claim: under
delayed-query drift, a cheap sub-100\,ms Phase~1 refinement can
recover part of the quality lost to aggressive compression without changing the
decode operating point.
That already turns idle time from wasted latency into a usable
inference-quality budget.
Bigger claims about Phase~2, full LongBench integration, and the complete
quality-throughput frontier remain scale-up work for the larger-GPU study.
The general principle---using scheduled idle compute for inference quality
improvement---extends beyond KV caches to other quality-speed tradeoffs in
neural inference.

% ---------------------------------------------------------------
%  Impact Statement
% ---------------------------------------------------------------
\section*{Impact Statement}

This paper presents work whose goal is to advance the field of Machine
Learning by improving the efficiency of large language model inference during
agentic deployments.
The proposed \textbf{IdleKV} method reduces computational waste and improves
quality-throughput tradeoffs, which could lead to more efficient and accessible
AI systems.
There are many potential societal consequences of our work, none which we feel
must be specifically highlighted here.

% ---------------------------------------------------------------
%  LLM/Agent Usage Disclosure
% ---------------------------------------------------------------
\section*{LLM/Agent Usage Disclosure}

All technical content, system design, experimental methodology, and analysis
are the sole work of the authors.
No LLM was used to generate experimental results, figures, or algorithmic
designs.

\bibliography{idlekv_references}
\bibliographystyle{plainnat}

\end{document}
